%! Choosing a Sample: Four Sampling Techniques — Unit 1, Lesson 1.3 — Homework
% PILOT SHAPE (2026-09-06): modeled on AP Statistics lesson 1.4 minus its AP Practice page.
% This is the lesson's SCORED individual practice and the LAST component in the packet.
% Students START IT IN CLASS in the last ten minutes of the period, alone, while the
% teacher circulates for the second formative read; they finish it at home. Packet pages are
% the default; a Desmos activity may replace them on a given day (decided at Close & Assign).
% THIRD CONTEXT: the notes teach on Riverbend HS, the Guided Practice runs on Cedar Ridge
% Apartments, and these items replan the Harbor Point pool study from 1.2 — recall is not the
% skill.
% TWO PAGES MAX (user decision 2026-09-06): two boxes — items 1-2 on page 1, items 3-6 and
% the spiralbox on page 2 — so no item breaks across a page.
% Item 5's two blanks are the deliberate CONTRAST PAIR (Plan B stratified vs. Plan D cluster on
% the same population); its write-lines are the target misconception head-on — a random draw
% of clusters that are alike inside. Item 1 is the spiral to 1.2. Item 6 is PS.DC.2b.
% THE DATA — Harbor Point Community Pool: 1,500 members numbered 0001-1500; 100 to survey.
%   Generator 0412 1587 0973 0412 0288 -> members 412, (1587 out of range), 973, (0412 repeat), 288.
%   Strata 600/525/375 at 1/15 -> 40/35/25. Systematic k = 1500/100 = 15, start 9 -> 24, 39, 54.
%   Cluster: 4 of 25 swim sessions x 25 = 100.
% PAGE LOCKSTEP: the key is generated from this file; every work block is byte-identical.
\documentclass[12pt]{article}
\usepackage{saar-article}
\usepackage{saar-boxes}
\newcommand{\TallMath}[1]{$\displaystyle #1\rule[-1.4em]{0pt}{3.2em}$}
% Word bank strip for every fill-in cluster (user request 2026-09-06). Defined per document;
% the key inherits it and prints the same strip, so heights match.
\newcommand{\wordbank}[1]{\par\vspace{2pt}\noindent\fcolorbox{goldacc}{goldbg}{\parbox{\dimexpr\linewidth-2\fboxsep-2\fboxrule\relax}{\small\textbf{Word bank:} #1}}\par\vspace{4pt}}

\begin{document}

\pageheader{Unit 1, Lesson 1.3}{Homework}
% NAMESTRIP: no \namedateperiod here. The student writes their name once, on the
% cover this component is stapled behind.

\vspace{0.12in}

% Authored identically in homework/ and homework_key/.
\begin{remindbox}
\small
\textbf{This is your graded homework.} It is scored out of 12 and due at the start of the
first class after two study halls. You will start it in the last minutes of the period, so ask about anything that stalls
you before you leave, then finish it on your own.
\end{remindbox}

\vspace{0.08in}

\boxguard[20]
\begin{scenariobox}[Harbor Point Community Pool --- Do It Again, Properly]{cerulean}
Last time, a staff member asked the $75$ members who arrived on a Tuesday evening. The board
has thrown that study out and asked you to plan a new one: $100$ of the pool's $1{,}500$
members.

\vspace{2pt}
What the office knows: the membership list is complete and numbered $0001$ to $1500$. There
are $600$ individual memberships, $525$ family memberships, and $375$ senior memberships. The
week has $25$ scheduled swim sessions, and about $25$ members attend each one.
\end{scenariobox}

\vspace{0.08in}

\boxguard
\begin{notesbox}{Four Plans for One Survey}
Every plan below surveys $100$ members --- the plans differ only in \emph{who gets chosen}.
\wordbank{all 1,500 members \;\textbullet\; the 75 asked Tuesday evening \;\textbullet\; convenience \;\textbullet\; yes --- in range and new \;\textbullet\; no --- out of range \;\textbullet\; no --- already chosen \;\textbullet\; member 973 \;\textbullet\; member 288 \;\textbullet\; none \;\textbullet\; 40 \;\textbullet\; 35 \;\textbullet\; 25 \;\textbullet\; 100 \;\textbullet\; 24 \;\textbullet\; 39 \;\textbullet\; 54 \;\textbullet\; stratified \;\textbullet\; cluster}
\begin{enumerate}[leftmargin=1.6em, itemsep=9pt, topsep=3pt]

  \item The Tuesday-evening study.
        \par\vspace{4pt}
        Population: \blank{4.2cm} \hfill Sample: \blank{5cm}
        \par\vspace{8pt}
        It was a \blank{2.6cm} sample. Name one kind of member who had \emph{no chance} of
        being chosen by it.
        \par\writeline

  \item \textbf{Plan A --- simple random sample.} You number the members $0001$ to $1500$
        and the generator returns {\ttfamily 0412 \; 1587 \; 0973 \; 0412 \; 0288}. Decide
        each one.

        \begin{center}
        \renewcommand{\arraystretch}{1.4}
        \begin{tabularx}{\linewidth}{@{} c X p{3cm} @{}}
        \toprule
        \textbf{Returned} & \textbf{Use it? Why or why not} & \textbf{Member chosen} \\
        \midrule
        0412 & yes --- in range and new & member 412 \\
        1587 & \blank{4.6cm} & \blank{2.6cm} \\
        0973 & \blank{4.6cm} & \blank{2.6cm} \\
        0412 & \blank{4.6cm} & \blank{2.6cm} \\
        0288 & \blank{4.6cm} & \blank{2.6cm} \\
        \bottomrule
        \end{tabularx}
        \end{center}

\end{enumerate}
\end{notesbox}

\boxguard[30]
\begin{notesbox}{Four Plans for One Survey, continued}
\begin{enumerate}[leftmargin=1.6em, itemsep=4pt, topsep=2pt, start=3]

  \item \textbf{Plan B --- stratified sample} by membership type: $100$ of $1{,}500$ is $1/15$,
        so take $1/15$ of each stratum. Show the senior row.

        \begin{center}
        \renewcommand{\arraystretch}{1.2}
        \begin{tabularx}{\linewidth}{@{} l c X @{}}
        \toprule
        \textbf{Membership} & \textbf{Members} & \textbf{Number to survey} \\
        \midrule
        Individual & 600 & \blank{3.8cm} \\
        Family     & 525 & \blank{3.8cm} \\
        Senior     & 375 & \blank{3.8cm} \\
        \midrule
        \textbf{Total} & \textbf{1,500} & \blank{3.8cm} \\
        \bottomrule
        \end{tabularx}
        \end{center}

        \begin{work}
          \dfrac{100}{1500} &= \dfrac{1}{15} \\
          \dfrac{1}{15} \times 375 &= 25 \text{ members}
        \end{work}

  \item \textbf{Plan C --- systematic sample} from the numbered list. Find $k$; the random
        start lands on member $9$; write the next three members chosen.

        \begin{work}
          k = 1500 \div 100 &= 15
        \end{work}

        Next three: \blank{1.4cm} \qquad \blank{1.4cm} \qquad \blank{1.4cm}

  \item \textbf{Plan D --- cluster sample.} Draw $4$ of the $25$ swim sessions at random and
        survey everyone there ($4 \times 25 = 100$). Plans B and D both use groups: Plan B is
        \blank{2cm} and Plan D is \blank{2cm}. The four sessions drawn are three lap-swim hours
        and one fitness hour --- a random draw, every count correct. Explain why the results
        still might not describe all $1{,}500$ members.
        \par\writelines{2}

  \item The board says: ``Whatever we do, seniors must be heard --- they pay dues too, and they
        swim at different hours.'' Which plan should they run, and why does it deliver that?
        \par\writelines{2}

\end{enumerate}
\end{notesbox}

\boxguard[5]   % the two-line spiralbox needs about five baselines; a larger guard pushes it to a third page
\begin{spiralbox}
\textbf{Next lesson: Bias in Samples and Surveys.} A perfectly chosen sample can still mislead
--- who was left off the frame, who refused to answer, how the question was worded. That is
\textbf{bias}.
\end{spiralbox}

\end{document}
